\documentclass[../../../include/open-logic-section]{subfiles}

\begin{document}

\olfileid{sth}{replacement}{intro}
\olsection{परिचय}

प्रतिस्थापन हा स्वयंसिद्धक-रूपबंधच $\ZF$ आणि~$\Z$ यांतील फरक ठरवतो. \crefrange{sth:ordinals::chap}{sth:spine::chap} या प्रकरणांमध्ये आपण त्याचा वापर केला. आता अखेर या प्रकरणात आपण विचारू: प्रतिस्थापनाचे समर्थन करता येते का?

हा प्रश्न नेमका करण्यासाठी प्रतिस्थापन खरेच बरेच \emph{सामर्थ्यवान} आहे, हे लक्षात घ्यावे लागेल. ते किती सामर्थ्यवान आहे याची कल्पना या प्रकरणात (आणि पुन्हा \olref[sth][card-arithmetic][fix]{sec} मध्ये) येईल. त्यावरून त्याच्या समर्थनाची खरोखर गरज आहे, हेही दिसेल.

समर्थनाचे दोन प्रकार पाहू. साधारणपणे, स्वयंसिद्धकाच्या फलितांवरून त्याचे समर्थन करण्याचा प्रयत्न म्हणजे \emph{बाह्य} समर्थन; तर संबंधित गणिती संकल्पनांतूनच स्वयंसिद्धकाला पाठबळ मिळते असे दाखवण्याचा प्रयत्न म्हणजे \emph{आंतरिक} समर्थन. याचा अर्थ या प्रकरणात अधिक स्पष्ट होईल; तरी ही चर्चा फार मोठ्या विषयाची केवळ सुरुवात आहे. अधिक माहितीसाठी विशेषतः \citeauthor{Maddy1988a} (\citeyear{Maddy1988a} आणि \citeyear{Maddy1988b}) पाहा.

\end{document}
